\documentclass[11pt]{article}

\usepackage[margin=1in]{geometry}
\usepackage[T1]{fontenc}
\usepackage[utf8]{inputenc}
\usepackage{lmodern}
\usepackage{amsmath}
\usepackage{booktabs}
\usepackage{enumitem}
\usepackage{parskip}
\usepackage{array}
\usepackage{tabularx}
\usepackage[hidelinks]{hyperref}

\newcolumntype{Y}{>{\raggedright\arraybackslash}X}
\newcolumntype{P}[1]{>{\raggedright\arraybackslash}p{#1}}
\newcommand{\code}[1]{\texttt{#1}}
\setlength{\emergencystretch}{2em}

\title{\textbf{bitsign MVP: Technical Specification} \\ \large v1.3 draft}
\author{Technical team}
\date{29 August 2026}

\begin{document}

\maketitle

\begin{abstract}
\noindent Implementation specification for the bitsign MVP. It defines the client,
service, inference, contribution, privacy, and release requirements. Numbers marked
\textbf{[default]} are launch defaults, server-configurable without a client release.
Requirement keywords: \textbf{must} / \textbf{should} / \textbf{may}.
\end{abstract}

\setcounter{tocdepth}{2}
\tableofcontents

\section{System Overview}

\begin{center}
\code{iOS app -> Cloudflare Worker (front door) -> API service (managed container)}\\
\code{API service -> Worker outbound adapters -> R2 / D1 / KV / Queues}\\
\code{API service -> inference service (CPU or GPU container runtime)}
\end{center}

Components:

\begin{enumerate}[nosep]
  \item \textbf{iOS app} (SwiftUI, iOS 17+): capture, review, upload, transcript
        display, and contribution flows. TestFlight only.
  \item \textbf{Front door}: Cloudflare Worker. TLS termination,
        streaming body and header limits, coarse IP rate limiting, request-id
        assignment, scheduled jobs, routing to the API container, and narrow
        outbound adapters for Cloudflare bindings.
  \item \textbf{API service}: Axum in a managed container. All business logic.
        It authenticates and authorizes requests. It is stateless; authoritative
        state is in D1, media is in R2, and KV contains non-sensitive configuration.
  \item \textbf{Inference service}: Axum with one deployment per runtime class:
        \code{cpu} and \code{gpu}. Reachable only from the API service through an
        authenticated private route; not exposed through the front door.
  \item \textbf{Shared domain library}: types, ids,
        error taxonomy, consensus scoring (pure functions, property-tested).
\end{enumerate}

All ids are ULIDs (26-char Crockford base32). All timestamps are UTC RFC 3339. All
API bodies are JSON with \code{snake\_case} keys.

\section{Accounts, Auth, and Devices}

\subsection{Sign-in and attestation}

Sign in with Apple only. Flow:

\begin{enumerate}[nosep]
  \item \code{POST /v1/auth/challenges} accepts the compile-time build channel,
        app version, and build number. The API host and an allowlist map them to the
        bundle id and App Attest environment. It returns a challenge id, 32 random
        bytes, Apple nonce, ten-minute expiry, and the expected environment. D1
        stores the binding and challenge hash and consumes the row atomically once.
  \item The client includes the nonce in Sign in with Apple. It hashes the returned
        identity token into the canonical App Attest client-data object, which also
        binds challenge, key, bundle, build channel, version/build number, and
        environment. A new key produces an attestation; a server-registered key
        produces an assertion.
  \item \code{POST /v1/auth/apple} accepts \code{identity\_token},
        \code{key\_id}, \code{kind}, either \code{attestation\_object} or
        \code{assertion\_object}, \code{client\_data}, and \code{challenge\_id}.
        The server verifies the single-use challenge; identity-token signature,
        issuer, configured client-id audience, expiry, and nonce; App Attest app id
        and environment; the attestation or assertion signature; and the monotonic
        assertion counter. It then upserts the account by Apple \code{sub},
        registers or advances the device key, and mints a session.
\end{enumerate}

\textbf{Session token}: opaque 32-byte random value returned once. D1 stores only
its SHA-256 with account id, device id, expiry, refresh, and revocation timestamps.
The initial lifetime is 30 days \textbf{[default]}. When fewer than 15 days remain, a
primary-D1 request may extend expiry to the earlier of 30 days from that request or 90
days from session creation, with at most one write per 24 hours. Client sends
\code{Authorization: Bearer}.
Revocation sets \code{revoked\_at} before the response returns. No JWTs and no KV
authorization cache.

Accounts have exactly one active attested device at MVP; signing in on a new device
supersedes the old one (old sessions revoked). Requests from builds that cannot
attest are rejected with \code{attestation\_required}. Debug simulator tests use an
injected provider against local configuration; distributed Staging and Production
builds have no bypass. App Attest keys can be lost on reinstall, restore, or device
migration, so an unknown key returns \code{attestation\_registration\_required} and
requires a fresh challenge and attestation. Assertions on every ordinary API request
are out of scope; returning sign-in uses an assertion.

The client stores a session record (token, bitsign account id, and Sign in with Apple
user id) and a separate App Attest key-id record in the Keychain using
\texttt{kSecAttrAccessibleAfterFirstUnlock}\allowbreak\texttt{ThisDeviceOnly}.
On \code{unauthorized}, the app checks Apple credential state using the stored Apple
user id, clears the session record, and presents Sign in with Apple when a new identity
token is required. It leaves the App Attest key-id record for a returning assertion and
does not claim silent token renewal.

\textbf{Deletion/recreation.} A deleted account's Apple \code{sub} may sign up
again and receives a fresh account. The new account does not restore deleted
content, history, or device registrations.

\subsection{Roles}

\code{accounts.role} $\in$ \{\code{user}, \code{internal}, \code{admin}\}.
Models scoped to internal evaluation (\S\ref{sec:models}) are servable only to
\code{internal} and \code{admin}. Named technical testers receive \code{internal};
external pilot participants remain \code{user}. Pilot access is a separate
authorization and never grants access to internal-evaluation models.

\section{Media Pipeline}

\subsection{Clip constraints}

\begin{center}
\begin{tabular}{@{}ll@{}}
\toprule
Property & Constraint \\
\midrule
Container / codec & MP4, H.264, exactly one video track and no audio track \\
Duration & 2.0--6.0\,s \\
Resolution & transformed long edge 480--1280, short edge 360--720 \\
Frame rate & nominal 30\,fps; server accepts 24--31 and at most 190 frames \\
Size & $\leq$ 12\,MiB \\
\bottomrule
\end{tabular}
\end{center}

The server \textbf{must} validate container, codec, duration, dimensions, and size
by demuxing (not by trusting metadata), and \textbf{must} reject violations with
\code{invalid\_media} and no partial processing. A versioned malformed/adversarial
fixture set runs in CI; every fixture must be rejected without panic, hang, or
residual object. Validation uses pinned FFmpeg/ffprobe 9.0.1, a full decode, exact
rational frame-rate arithmetic, and the metadata and orientation policy in
\code{IMPLEMENTATION\_PLAN.md}.

\subsection{Upload protocol}

\begin{enumerate}[nosep]
  \item {\small\sloppy Upload reservation uses \code{POST /v1/uploads}. Its JSON body contains
        \code{purpose}, \code{content\_length}, and \code{sha256}. Purpose is
        \code{communicate}, \code{task\_submission}, or, after E9,
        \code{pilot\_evaluation}. The response carries \code{upload\_id},
        \code{put\_url}, \code{required\_headers}, and \code{expires\_at}.
        A Communicate PUT expires after ten minutes; a
        contribution or pilot PUT expires after 60 minutes. The client sends exactly
        the signed headers. Only contribution uses a background task; it starts only
        when at least 15 minutes remain and obtains a new reservation after expiry.
        The server-generated object key is the purpose, a slash, and a random ULID. It
        contains no account, task, script, participant, filename, or model identifier.\par}
  \item Client PUTs the clip, then references \code{upload\_id} in the follow-up
        call. Before any use, the server verifies object existence and that the
        object's actual size and SHA-256 match the declaration; mismatch $\to$
        \code{invalid\_media} and immediate object deletion.
  \item Communicate, Contribute, and pilot evidence use separate storage buckets.
        The first is ephemeral; the others follow their respective consent and
        retention policies. \code{purpose} selects the bucket at presign time; an
        upload can never move between them.
\end{enumerate}

\textbf{Ephemerality mechanism} (Communicate): (a) inline delete after successful
inference; (b) an hourly Worker job deletes any ephemeral communication object
older than 24\,h and marks its \code{uploads} row \code{deleted}; (c) a bucket
lifecycle rule at the minimum granularity R2 supports (age $\geq$ 1 day) as the
final backstop. The guarantee that is tested and stated publicly is (a)+(b); the
lifecycle rule is defense in depth, not the mechanism.

\section{Inference}
\label{sec:infer}

\subsection{API}

\code{POST /v1/inferences} accepts \code{upload\_id}, \code{mode}
(continuous or bounded), and the canonical \code{config\_version} in its body,
requires the \code{Idempotency-Key} header, and returns a synchronous response.

The response carries \code{request\_id} and \code{transcript}; a model object with
\code{id}, \code{mode}, \code{version}, \code{capability\_label}, and
\code{release\_scope}; and timings with \code{queue\_ms} and
\code{inference\_ms}.

Rules: the idempotency record is written before dispatch, so a retried request
cannot duplicate processing. It stores no transcript. Server wall-clock timeout 30\,s $\to$
\code{inference\_timeout}; the client \textbf{must} use a request timeout of 35\,s.
After a transport failure, the client may retry with the same idempotency key to
learn whether the operation is still pending. A completed operation whose response
was lost returns \code{result\_unavailable}; the app creates a new upload from the
still-local clip when available, or offers recapture after that clip is gone. The
server does not retain the transcript. One in-flight inference per account
(\code{conflict} otherwise). The API service resolves \code{mode} $\to$ active
model through current authenticated configuration and the registry, verifies the
echoed mode and RFC-8785 config digest before dispatch, and routes to the recorded
runtime. A mismatch returns \code{config\_stale}; client fields are not authority. Transcripts are
returned and never stored server-side (\S\ref{sec:schema}).

Precise semantics an implementer must not invent:

\begin{itemize}[nosep]
  \item \textbf{Idempotent retry}: same \code{idempotency\_key} while the
        original is in flight $\to$ \code{conflict}; after successful completion
        $\to$ \code{result\_unavailable}. D1 retains terminal metadata and the
        request hash for 24\,h, never response bytes, transcript text, or an
        output-derived hash.
  \item \textbf{\code{queue\_ms}} is time spent waiting for an inference
        concurrency slot (\S\ref{sec:inferservice}); \code{inference\_ms} is model
        execution only.
  \item \textbf{Upload reuse}: a \code{stored} upload may be retried with a new
        idempotency key after a failed inference; it becomes \code{consumed}
        (comm object deleted inline) only on success. The hourly sweep remains
        the failure backstop.
  \item \textbf{Cold starts}: if the routed runtime reports not-ready
        (\S\ref{sec:inferservice}), the API returns \code{model\_warming} with
        \code{retry\_after\_ms}; the client shows a warming state and retries
        automatically up to 3 times. Warming is expected with scale-to-zero GPU
        and is not an error condition in telemetry.
\end{itemize}

\subsection{Model registry and serving policy}
\label{sec:models}

D1 table \code{models}: \{\code{id}, \code{mode}, \code{version},
\code{release\_scope}: \code{internal\_evaluation} $|$ \code{approved\_pilot},
\code{runtime}: \code{cpu} $|$ \code{gpu}, \code{endpoint}, \code{weights\_sha256},
\code{capability\_label}: \code{technical\_prototype} $|$ \code{bounded\_vocabulary}
$|$ \code{short\_clip\_preview}, \code{active}\}. Exactly one active row per
\code{mode} (partial unique index). The API \textbf{must} refuse to serve a
model scoped to \code{internal\_evaluation} to a non-internal account
(\code{model\_unavailable}), regardless of registry state. This rule is enforced
in the request path and covered by CI.

The checked-in fixture registry maps rights-cleared fixture hashes to fixed outputs,
has internal-evaluation release scope, and can prove only the connection
path. It returns \code{model\_unavailable} for an unknown hash and is rejected by
Production configuration. A real model is added only from a signed-off manifest that
records source, revision, weights hash, rights determination, preprocessing, decoding,
hardware, and approved release scope. A bounded model also requires an approved
top-50 curriculum and at least 20 accepted clips per class \textbf{[default]} before
evaluation; neither model nor curriculum is selected by an implementation agent.

\subsection{Inference service interface}

\label{sec:inferservice}
\code{POST /internal/inferences} streams a \code{video/mp4} body with
\code{Content-Length}, \code{X-Video-SHA256}, \code{X-Model-ID}, and
\code{X-Request-ID}, authenticated with a distinct internal credential. The runtime
verifies the 12\,MiB streaming bound, length, and digest before returning
\{\code{transcript}, \code{inference\_ms}, optional internal diagnostics\}. No R2
object key, presigned URL, account or task id, prompt, or capability decision crosses
this boundary. The contract is language-neutral: the
fixture runtime is Rust, while an approved real model may use Rust or Python. The
service loads the model named by its manifest at boot and fails fast on a weights-hash
mismatch. Limits: max decoded frames 190; output is UTF-8 plain text with at most
2 KiB and 120 Unicode-whitespace-delimited words after trimming and whitespace
collapse; max wall 25\,s; max concurrency 2 \textbf{[default]} with a
bounded wait queue of 4 \textbf{[default]} (max slot wait 10\,s, then
\code{overloaded}); typed errors; never fabricated output. Undecodable media or
an unloaded model is an error, not a placeholder string.

The public API and app expose no numeric confidence. The app renders validated output
as plain text without Markdown, HTML, or link detection. Invalid or over-limit output is an error and is never truncated
into a plausible translation. Pilot evidence stores exactly the validated string
returned to the app.

Readiness: \code{GET /internal/ready} returns \code{ready} only when weights are
loaded and a self-test inference on a bundled 1\,s fixture clip has succeeded since
boot. The API's pre-dispatch check uses this endpoint (cached 5\,s) to produce
\code{model\_warming}; demo pre-warm automation polls it to completion.

\section{Contribute Pipeline}
\label{sec:contribute}

\subsection{Consent}

First entry into Contribute presents the study consent (purpose: model training and
evaluation; retention; access roles; withdrawal and deletion; explicit statement
that Communicate clips are never used). Consent is versioned:
\code{consents(account\_id, version, granted\_at, withdrawn\_at)}. Tasks are issued
only to accounts whose consent version is current; a bumped version re-prompts.
Withdrawal stops issuance immediately and triggers \S\ref{sec:deletion} handling
for that account's contributed raw and derived material.

MVP contribution accounts are named adults age 18 or older. The app does not request
or infer age from Sign in with Apple; eligibility is established during controlled
enrollment. Supporting minors requires a later consent and product protocol.

\subsection{Scripts and tasks}

\code{scripts}: \{\code{id}, \code{ciphertext} (AES-256-GCM; key held as a
container secret, distinct from session material), \code{sha256} (of plaintext),
\code{class}: \code{fingerspell} $|$ \code{sign} $|$ \code{phrase}, \code{active}\}.
Script plaintext appears in exactly one API response type (Task~A issuance) and is
never logged.

\textbf{Issuance policy.} \code{GET /v1/tasks/next} returns, in priority order: an
open Task~B review if one is assigned (reviews gate pipeline latency), else a new
Task~B assignment if the account qualifies, else a new Task~A. At most one open
item of \emph{each} type per account. Scripts and aliases are admin-seeded
(\S\ref{sec:admin}); the launch set is the fingerspelling alphabet, digits, and
the approved top-50 sign curriculum.

\textbf{Task A (meaning $\to$ sign).} Issues \{\code{task\_id},
\code{script\_text}, \code{class}, \code{expires}\} (24\,h expiry).
An account is never issued the same script twice (unique
\code{(script\_id, account\_id)}). Submission:
\code{POST /v1/tasks/\{id\}/submit} \{\code{upload\_id}\}; the server records
\code{video\_sha256} and enqueues review assignment.

\textbf{Task B (sign $\to$ meaning).} The assignment job selects $k=3$
\textbf{[default]} reviewers: consented, not the contributor, not the contributor's
device, and not sharing the contributor's network neighborhood, checked via
\code{net\_hash} = HMAC(salt, /24 or /48 prefix) stored on submissions and reviews
and purged after 30 days \textbf{[default]}; the raw address is never stored.
Selection is weighted toward higher reviewer accuracy. Reviewers receive
\{\code{review\_id}, presigned GET (1\,h)\}, never the script. Reviews expire
and reassign after 24\,h. Reviewers discover work by polling \code{GET
/v1/tasks/next} (which returns either task type); there are no push notifications
at MVP.

\subsection{Consensus scoring}
\label{sec:consensus}

Pure, deterministic, versioned function (\code{scoring\_v1}), stamped on every
result. Normalization: Unicode NFKC, lowercase, strip punctuation, collapse
whitespace.

Per-review match score against the script's \emph{reference set}
$R = \{\text{script}\} \cup \{\text{accepted references}\}$, seeded with the
script itself, so scoring works before any reference exists:
\[
m_i \;=\; \max_{r \in R} \; \max\!\bigl(0,\; 1 - \mathrm{WER}(\text{recon}_i,\, r)\bigr)
\]
For \code{fingerspell} and \code{sign} classes, $m_i \in \{0,1\}$ by exact
normalized match against $R$ extended by the script's alias table. Consensus
$C = \mathrm{median}(m_1, \ldots, m_k)$.

Outcomes \textbf{[defaults]}: $C \geq 0.85$ \code{accepted}; $0.70 \leq C < 0.85$
\code{variant} (admin secondary review); $C < 0.70$ \code{flagged} (no contributor
acceptance pending investigation; clip retained as edge data unless invalid).

\textbf{Reference-set growth}: an accepted clip's reconstructions with
$m_i \geq 0.85$ are candidates for $R$; admission requires normalized edit distance
$\geq 0.15$ from every existing member (no near-duplicates), cap $|R| \leq 5$
\textbf{[default]} excluding the script. Growth applies to \emph{future} scoring
only; completed results are never rescored.

\textbf{Quorum and timing}: consensus runs when all $k$ reviews are submitted, or
at the 24\,h review deadline with at least 2 submissions ($C$ = median of what
exists). Fewer than 2 at deadline $\to$ expired reviews reassign to fresh
reviewers; after 3 \textbf{[default]} reassignment cycles the task is escalated to
the admin queue rather than expiring against the contributor, whose submission
remains pending until an eventual decision.

\textbf{Anchor tasks}: 5\% \textbf{[default]} of issued reviews reference clips
with known-correct reconstructions, indistinguishable client-side. Anchor outcomes
update \code{reviewer\_stats.accuracy} (EWMA, $\alpha = 0.2$); accuracy $< 0.6$
\textbf{[default]} sets \code{hold}; held accounts receive no new contribution
assignments pending admin review. Cold start: the launch anchor set is team-recorded
clips against known scripts; thereafter, \code{accepted} clips with unanimous
$m_i \geq 0.85$ are eligible for promotion to anchors (admin action), which is also
how the anchor pool grows with the corpus.

\subsection{Admin surface}
\label{sec:admin}

\code{role=admin} only; served by the same API and exercised by CLI/HTTP at MVP.
There is no admin UI. Operations are:

\begin{itemize}[nosep]
  \item \code{GET /v1/admin/queue} and \code{GET /v1/admin/metrics};
  \item \code{POST /v1/admin/decisions} and \code{POST /v1/admin/scripts};
  \item \code{POST /v1/admin/holds/release}, with the account id in the body; and
  \item \code{POST /v1/admin/anchors/\{task\_id\}/promote}.
\end{itemize}

These cover variants, flags, escalation, holds, script and alias management, anchor
promotion, queue depth, consensus latency, and reviewer accuracy. Every admin action
writes an audit row.

\section{Deletion and Withdrawal}
\label{sec:deletion}

Before \code{DELETE /v1/me}, the app generates and Keychain-stores a random 32-byte
status secret. The request sends it once as \code{Deletion-Status-Token}; the server
stores only its SHA-256, returns the job id and initial status, and queues a deletion
job (target: minutes; completion within 30 days \textbf{[default]}): revoke sessions; delete R2
objects in every purpose bucket belonging to the account plus derived artifacts; delete
account PII and device rows; remove the account's reconstructions from reference
sets; void in-flight consensus involving the deleted material; purge
\code{net\_hash} values.
\code{GET /v1/deletions/status} reports status without account authentication because
the account may already be gone. It accepts the bearer secret only as
\code{Authorization: Deletion <token>}; the token never appears in the URL and only
its hash is stored. A lost deletion response does not lose status access because the
app persisted the secret before the request. The token permits repeated status reads,
expires 45 days after job creation, and is then removed. The job emits a
minimal completion record containing the job id and timestamps, with no account
identifier or contributed content.

\section{iOS Application}

\subsection{Screens and required behavior}

\begin{center}
\begin{tabularx}{\textwidth}{@{}P{3.2cm}Y@{}}
\toprule
Screen & Must-behavior \\
\midrule
Onboarding & Capability-label explanation, ``aid, not interpreter replacement''
framing, and explanation of permissions. Camera, microphone, and Speech permissions
are requested only when their feature is invoked. Blocks until Sign in with Apple and
attestation complete. \\
Communicate & Camera preview with framing guide (face + torso + both hands); record
button with 6\,s hard stop and countdown ring; clip review (play, retry, discard,
send); rolling transcript list, \textbf{in-memory only, cleared on app
termination; nothing is persisted on device or server}, with each entry showing text,
capability-label badge, speak button, and an optional local text edit. Ordinary
Communicate has no rating control or feedback endpoint; only a separate pilot or
research protocol may submit ratings or corrections. Reply view: mic button $\to$ on-device
\code{SFSpeechRecognizer}, when supported, $\to$ large-type text, with typed fallback.
Upload and inference progress are
visually distinct; every error state names a recovery action, and recoverable upload
or inference failures preserve the recorded clip for retry. \\
Contribute & Consent flow (scrollable, versioned, decline path exits cleanly); task
screen showing the one prioritized next item; Task A capture reusing the Communicate
capture component; Task B player + reconstruction field; protected submissions queue
with the limits, expiry, reconciliation, and retry states in
\code{IMPLEMENTATION\_PLAN.md}. \\
Settings & App/model versions and capability label, consent status + withdrawal,
account deletion (typed confirmation, then status screen), notices. \\
\bottomrule
\end{tabularx}
\end{center}

E9 adds a server-authorized Pilot task surface and a separate Review surface visible
only while \code{pilot\_reviewer} permission is unexpired. Pilot transfer is foreground
only and has no offline queue. These surfaces follow the local-file and evidence
lifecycles in \code{IMPLEMENTATION\_PLAN.md}.

\subsection{Capture component}

\code{AVCaptureSession} at the 720p preset, 30\,fps, front camera; MP4/H.264 with no
audio input or track; segments 2--6\,s; accepts portrait or landscape through the
track transform; writes ordinary clips to a protected app-container cache and queued
contributions to protected Application Support. Files are excluded from backup,
deleted according to their flow, and never enter the photo library. A recording
indicator is always visible while capturing. The live front-camera preview is
mirrored for the signer, while the encoded file is canonical and unmirrored; a fixture
verifies handedness and one-time orientation-transform application.

\subsection{Accessibility (release-gating)}

VoiceOver labels and logical focus order on every screen; Dynamic Type through
\code{accessibility5}; contrast $\geq$ WCAG AA; every audio cue has a visual
equivalent; the speak button is optional, never required to read a result.

\section{Front Door and Cross-Cutting Requirements}
\label{sec:frontdoor}

\begin{itemize}[nosep]
  \item \textbf{Request ids}: the Worker assigns a ULID \code{request\_id}, echoed
        in every response and log line.
  \item \textbf{Wire bounds}: request target $\leq 2$\,KiB, decoded headers
        $\leq 16$\,KiB, and public JSON request and response bodies $\leq 64$\,KiB,
        except admin script import at $\leq 256$\,KiB. The Worker and API enforce
        the applicable decoded bound. Direct media remains capped at 12\,MiB and is
        not proxied through a public JSON route.
  \item \textbf{Rate limits}: after authentication, D1-backed atomic token buckets
        are keyed by session. Reads have capacity 240 and refill at 2 tokens/s;
        mutations have capacity 60 and refill at 0.5 tokens/s. A request costs one
        token. Integer millitoken arithmetic and server time are mandatory. The
        Worker also applies a coarse IP backstop and limits each auth route to
        10 requests/min/IP without persisting or logging the address. Exhaustion
        returns \code{429}, \code{Retry-After}, and \code{retry\_after\_ms}.
  \item \textbf{Error envelope}: non-200 responses carry required \code{code},
        \code{message}, and \code{request\_id} fields under \code{error}, plus
        optional \code{retry\_after\_ms}. Stable codes are
        \code{unauthorized}, \code{attestation\_required},
        \code{attestation\_registration\_required}, \code{consent\_required},
        \code{invalid\_media}, \code{media\_too\_long}, and
        \code{upload\_expired}; \code{task\_expired}, \code{config\_stale}, and
        \code{result\_unavailable}; \code{model\_unavailable},
        \code{model\_warming}, and \code{inference\_timeout}; plus
        \code{rate\_limited}, \code{overloaded}, \code{not\_found},
        \code{conflict}, and \code{internal}. Clients switch on code, never on
        message.
  \item \textbf{Telemetry}: logs and metrics carry request id, build and model
        version, sizes, timings, status, and error code. A daily-rotated HMAC
        pseudonym may support abuse correlation. They \textbf{must never} carry
        account id, video bytes, frames, transcripts, output hashes,
        reconstructions, script text, Apple identifiers, presigned URLs, or raw
        client addresses. Contribution \code{net\_hash} stays in D1 only where
        \S\ref{sec:contribute} requires it.
  \item \textbf{Secrets}: no static secret ships in the app; server secrets live in
        Worker and container secret bindings; the script-encryption key, the
        \code{net\_hash} salt, and the inference-service bearer are distinct
        secrets, rotatable independently.
  \item \textbf{Config}: every \textbf{[default]}-marked value lives in KV under
        \code{cfg:*}, hot-reloaded by the API service on a 60\,s cycle; config
        reads never block the request path. KV contains no session, consent,
        deletion, idempotency, or role authority.
  \item \textbf{Cron} (Worker scheduled events): hourly comm-bucket sweep
        (\S3.2); daily expiry of stale tasks/reviews; daily \code{net\_hash} purge;
        daily deletion-job retry.
\end{itemize}

\section{Operations}

\begin{itemize}[nosep]
  \item \textbf{Environments}: \code{staging} and \code{prod}, fully isolated with
        separate Worker deployments, D1 databases, KV namespaces, R2 buckets,
        inference deployments, and secrets. TestFlight builds point at \code{prod};
        internal TestFlight builds may point at staging, while external pilot builds
        point at production. Staging may use fixtures and shortened scheduled-job
        intervals.
        Environment is baked into the build; no runtime switcher exists in the app.
  \item \textbf{Migrations}: D1 schema changes are numbered and applied through
        the deployment tool. CI fails if a model or type change lands without a
        migration.
  \item \textbf{Health}: \code{GET /healthz} is shallow (process up, config
        loaded) and public. \code{GET /readyz} is deep (D1, KV, R2 reachable;
        active models' readiness per \S\ref{sec:inferservice}) and restricted; the
        pre-warm script polls \code{/readyz} until green.
  \item \textbf{CI gates}: \code{cargo fmt --check}, \code{clippy -D warnings},
        unit + property tests, the malformed-media fixture suite, and a
        staging deploy smoke test (\code{/readyz} green + one fixture inference)
        before any prod deploy.
  \item \textbf{Deploys}: Worker and Cloudflare containers use a pinned Wrangler
        release. An external GPU runtime, if an approved model requires one, uses a
        pinned provider command recorded with the image digest. Rollback redeploys
        the previous digest and restores the prior registry row; both operations
        must be scripted.
\end{itemize}

\section{Data Model (D1)}
\label{sec:schema}

Tables (PK first; FKs by name; all with \code{created\_at}):

{\small\sloppy
\begin{itemize}[nosep]
  \item \textbf{accounts}: \code{id}, \code{apple\_sub} (unique), \code{role},
        and \code{status} (active, hold, or deleted).
  \item \textbf{auth challenges}: \code{id}, \code{challenge\_sha256},
        \code{apple\_nonce}, \code{expected\_bundle},
        \code{expected\_environment}, \code{expires\_at}, and nullable
        \code{consumed\_at}.
  \item \textbf{devices}: \code{id}, \code{account\_id},
        \code{attest\_key\_id}, \code{attest\_public\_key},
        \code{attest\_counter}, \code{attest\_environment}, and nullable
        \code{superseded\_at}.
  \item \textbf{sessions}: primary key \code{token\_sha256},
        \code{account\_id}, \code{device\_id}, \code{expires\_at}, and nullable
        \code{last\_refreshed\_at} and \code{revoked\_at}.
  \item \textbf{rate limit buckets}: composite key \code{(token\_sha256, class)},
        \code{millitokens}, and \code{refilled\_at}; rows expire 24 hours after
        their session expires.
  \item \textbf{consents}: \code{account\_id}, \code{version},
        \code{granted\_at}, and nullable \code{withdrawn\_at}.
  \item \textbf{uploads}: \code{id}, \code{account\_id}, \code{purpose},
        \code{bucket}, \code{object\_key}, \code{declared\_sha256},
        \code{declared\_bytes}, nullable \code{verified\_bytes}, and \code{status}.
        Purpose is communicate, task submission, or pilot evaluation. Status is
        pending, stored, consumed, or deleted.
  \item \textbf{inferences}: \code{id}, \code{account\_id}, \code{upload\_id},
        \code{model\_id}, \code{status}, nullable \code{error\_code}, and
        \code{inference\_ms}. There is no transcript column.
  \item \textbf{idempotency records}: \code{account\_id}, \code{operation},
        \code{key}, \code{request\_sha256}, \code{status}, nullable
        \code{terminal\_code}, and \code{expires\_at}. There are no response bytes,
        transcript text, or output-derived hashes.
  \item \textbf{scripts}: \code{id}, \code{ciphertext}, \code{sha256},
        \code{class}, and \code{active}.
  \item \textbf{script aliases}: \code{script\_id} and
        \code{normalized\_alias}.
  \item \textbf{script references}: \code{script\_id}, \code{text},
        \code{source\_review\_id}, and \code{rank}; at most five rows per script.
  \item \textbf{tasks}: \code{id}, \code{script\_id}, \code{account\_id},
        \code{type}, \code{status}, nullable \code{upload\_id},
        \code{video\_sha256}, and \code{net\_hash}, plus \code{expires\_at}.
        Status is issued, submitted, or expired. The script and account pair is
        unique.
  \item \textbf{reviews}: \code{id}, \code{task\_id},
        \code{reviewer\_account\_id}, \code{status}, nullable
        \code{reconstruction}, \code{match\_score}, and \code{net\_hash}, plus
        \code{is\_anchor} and \code{expires\_at}. Status is issued, submitted, or
        expired.
  \item \textbf{consensus results}: primary key \code{task\_id},
        \code{scoring\_version}, \code{score}, \code{outcome}, and
        \code{decided\_at}. Outcome is accepted, variant, or flagged.
  \item \textbf{reviewer stats}: primary key \code{account\_id},
        \code{accuracy\_ewma}, \code{reviews\_completed}, and
        \code{anchors\_seen}.
  \item \textbf{models}: \code{id}, \code{mode}, \code{version},
        \code{release\_scope}, \code{runtime}, \code{endpoint},
        \code{weights\_sha256}, \code{capability\_label}, and \code{active}.
        Active rows are unique by mode.
  \item \textbf{deletion jobs}: \code{id}, nullable \code{account\_id} with
        \code{ON DELETE SET NULL}, \code{status},
        \code{status\_token\_sha256}, and nullable \code{completed\_at}. The account
        link is cleared when its row is purged; the status hash remains only through
        its 45-day expiry.
\end{itemize}
}

E9 adds separate pilot-consent, task, exact-output, rating, review, access-audit, and
withdrawal tables. Those tables may reference a \code{pilot\_evaluation} upload but
ordinary \code{inferences} and \code{communicate} uploads cannot be promoted into
them.

\code{net\_hash} columns are purged (set NULL) by the daily cron after 30 days
\textbf{[default]}.

\section{Non-Functional Requirements}

\begin{center}
\begin{tabular}{@{}lll@{}}
\toprule
Metric & Target & Measured as \\
\midrule
End-to-end translate & p95 $\leq$ 12\,s warm & tap Translate $\to$ rendered text \\
Upload leg & p95 $\leq$ 5\,s on LTE profile & PUT complete \\
Inference leg & p95 $\leq$ 6\,s warm & \code{inference\_ms} \\
API availability & best effort; pre-warm before demos & \code{/healthz} probes \\
Crash-free flow & 20 consecutive scripted runs & UI test suite \\
Reliability suite & $\geq$ 95/100 warm successes & 10 fixed clips $\times$ 10 \\
\bottomrule
\end{tabular}
\end{center}

Cold starts are acceptable at MVP; any scheduled demo or session is preceded by a
pre-warm operation (one inference per active model).

\section{Acceptance Tests}
\label{sec:accept}

Demo gates:

\begin{enumerate}[nosep]
  \item \textbf{D1}: TestFlight build installable by an authorized tester on a
        physical iPhone.
  \item \textbf{D2}: \code{/healthz} public at a documented URL; a valid
        \code{/v1/inferences} request returns a transcript from the declared model.
  \item \textbf{D3}: recorded session: signing captured in the TestFlight build
        returns an inference result originating from the live backend.
\end{enumerate}

Engineering gates (must pass before D3 is recorded):

\begin{enumerate}[nosep]
  \item Malformed-media fixtures: 100\% rejected, zero panics, zero residual R2
        objects.
  \item Idempotency: concurrent retries cause at most one inference execution;
        completed ordinary inference returns \code{result\_unavailable} without a
        stored transcript or second dispatch.
  \item Attestation: replayed or expired challenges are rejected. An unattested
        build receives \code{attestation\_required}.
  \item Serving policy: a model scoped to \code{internal\_evaluation} requested by
        a \code{role=user} account $\to$ \code{model\_unavailable}, under every
        registry state.
  \item Ephemerality: successful inference deletes the comm object inline; a
        deliberately failed inference's object is gone after the next hourly
        sweep; \code{uploads.status} transitions verified.
  \item Deletion: server integration tests seed ordinary account, device, session,
        upload, and inference records; call \code{DELETE /v1/me}; verify the R2 and
        D1 purge plus dead presigned URLs; and verify the status endpoint. Separate
        iOS tests verify Keychain, temporary-file, and in-memory cleanup. A server
        test never claims to erase an iPhone file.
  \item Warming path: cold GPU runtime $\to$ \code{model\_warming} with
        \code{retry\_after\_ms}; subsequent retry succeeds; \code{/readyz}
        reflects both states.
  \item Reliability + latency suite per the NFR table, recorded with device,
        network profile, and percentile method.
  \item Accessibility audit checklist on every screen present in the internal build.
\end{enumerate}

Before the Contribution build, its additional engineering gates must pass:
consensus is deterministic, reviewer-order-invariant, and version-stamped;
reference-set growth never changes a completed result; a deadline with 2 of 3
reviews scores their median, 1 review reassigns, and 3 reassignment cycles escalate;
deletion purges contribution clips, reviews, queued files, and derived records; and
the Contribute accessibility checklist passes.

\section{Build Plan}

The executable sequence, dependency graph, repository layout, commands, and exit tests
are in \code{IMPLEMENTATION\_PLAN.md}. Its work packages E0--E9 and I0--I5 separate
the internal vertical slice, contribution build, and closed pilot. An implementation
change that alters this specification must update the execution plan and OpenAPI in
the same commit.

\section{Out of Scope (MVP)}

Streaming inference; reverse sign generation; avatars; multilingual routing;
Android; web; App Store public release; per-request App Attest assertions; push
notifications; distributed-ledger integration; enterprise features; natural-voice
cloning; production HA. Pilot-mode implementation begins through E9 only after R1,
R2, and its protocol are approved. External recruitment and sessions remain prohibited
until the resulting candidate build passes R3--R5 as well.

\end{document}
